
\begin{remark}{7.2.303}\label{rem:7-2-303}
\par\noindent\textup{(1)}
我们已经证明，
$\mathfrak f^{\operatorname{id}\times\phi'}$
等滤过 $A_{n,K}$ 同态与
$\operatorname{Eval}$ 相容：
严格等式 (7.2.301.1)、(7.2.301.2) 成立。
但我们并未断言它们与
$\operatorname{Incl}$ 相容。
它们与 $\operatorname{Eval}$ 的相容性，
来自构造
$\mathfrak f^{\operatorname{id}\times\phi'}$
等映射所用模空间的多值截面
必须满足的相容条件。
另一方面，要定义
$\mathfrak f^{\operatorname{id}\times\phi'}$
使下式严格成立，似乎很困难：\translatornote{原书右端写
$\operatorname{Incl}_s\times\mathfrak f^{\phi'}$，
而本注 (2) 的相容等式用
$\operatorname{Incl}\circ\mathfrak f^{\phi'}$。
这里保留源书的乘号，待后文构造复核。}
\[
\mathfrak f^{\operatorname{id}\times\phi'}
\circ\operatorname{Incl}_s
=\operatorname{Incl}_s\times\mathfrak f^{\phi'}.
\]
原因在于，若用模空间 (7.2.302)
定义 $\mathfrak f^{\operatorname{id}\times\phi'}$，
则当 $(k,\beta)\neq0$ 时，
$(\operatorname{Incl}_s)_{k,\beta}\neq0$。
这一点与第 7.2.12 小节的
命题~\ref{prop:7-2-217}，
以及后文将证明的
命题~\ref{prop:7-2-277} 有关。

\par\noindent\textup{(2)}
实际上，也可利用
定理~\ref{thm:7-2-176}
定义映射\translatornote{原书此处将靶空间的底
写为 $L'_0$；引理 7.2.300 的对应映射
及下文包含相容式均以 $L'_1$ 为靶，
这里据源靶型别补正。}
\[
\operatorname{id}\times\mathfrak f^{\phi'}\colon
C([0,1]_s\times L'_0,\lambda'_0)
\longrightarrow
C([0,1]_s\times L'_1,(\lambda'_0,\lambda'_1)).
\]
它满足 (7.2.301)，还满足
\[
(\operatorname{id}\times\mathfrak f^{\phi'})
\circ\operatorname{Incl}
=\operatorname{Incl}\circ\mathfrak f^{\phi'}.
\]
换言之，在证明引理~\ref{lem:7-2-300}
的范围内，可以用纯代数方法定义这个映射，
使它扮演
$\mathfrak f^{\operatorname{id}\times\phi'}$
的角色。

不过，在引理 7.2.327 的证明中，
我们将使用
$\mathfrak f^{\operatorname{id}\times\phi'}$
的几何构造。
特别是要用到这样一个事实：
该映射由与模空间 (7.2.302)
的虚基本链相应的对应所定义。
见注 7.2.329。
因此，我们不采用上述
$\operatorname{id}\times\mathfrak f^{\phi'}$
映射。
\end{remark}
